\documentclass{wlwtxjsccp}

\begin{document}
\makecover
\tableofcontents

\begin{reportabstract}{通信协议、数据帧、热水站（3-5个）}
本报告设计了一套用于校园热水站智能监控系统的专用通信协议，旨在解决大规模、分布式热水设备的数据采集、状态监控与集中管理问题。协议面向低功耗、窄带宽的物联网场景（如4G/NB-IoT），采用主从式长连接架构，定义了精简高效的二进制报文格式、完备的通信流程及可靠的安全机制。
\end{reportabstract}

\cqutBeginReportBody

\section{项目概述}
随着高校后勤管理数字化、智能化转型的深入，传统的校园热水供应系统（如澡堂、开水房、宿舍淋浴间）面临诸多挑战：计费不透明、故障响应慢、数据统计难、能源浪费大、用户体验差。为解决这些问题，建设基于物联网（IoT）的智能热水供应系统已成为必然趋势。

本报告旨在定义该系统中最为核心的组成部分——\textbf{智能水表（终端）}与\textbf{集中器（云平台之间）}的通信协议。该协议是确保数据准确采集、指令可靠下发、系统稳定运行的关键，是实现智能计量、远程管控、数据分析的基础。

\section{需求分析}
\subsection{用户需求分析}
用户需求主要来源于系统最终使用者（学生）、运营管理者（后勤部门）和维护人员。

\blocklabel{A. 学生（最终用户）需求：}
\leaditem{透明计费}{实时、准确地显示用水量与消费金额，消费记录可便捷查询，杜绝“糊涂账”。}
\leaditem{稳定可靠}{刷卡或扫码后迅速出水，用水过程中无中断，结束响应及时。}
\leaditem{便捷支付}{支持校园一卡通、微信/支付宝扫码、APP充值等多种支付方式。}
\leaditem{故障反馈}{遇到设备故障（如不出水、关不停）时，能有简便渠道快速报修。}
\leaditem{节能省费}{可自主查询用水习惯，获取节能建议。}

\blocklabel{B. 后勤运营管理需求：}
\leaditem{精准计量与收费}{实现按实际用水量计费，杜绝包月制下的浪费，提升经济效益。}
\leaditem{实时监控}{远程实时监控所有热水终端的在线状态、工作状态（空闲、使用中、故障）、累计用水量等。}
\leaditem{远程控制}{可对指定终端进行远程开关阀、费率调整、设备重启等操作。}
\leaditem{财务与数据分析}{自动生成各类财务报表（日/月/年报表、单设备报表）、用水量分析（高峰时段、人均用水）、能耗分析，为决策提供支持。}
\leaditem{故障预警与告警}{设备离线、阀门异常、通讯中断、低余额、传感器故障等能主动上报，并定位到具体设备，实现预测性维护。}
\leaditem{参数统一下发}{能够批量对终端进行费率、时钟、开关机时间等参数同步。}

\blocklabel{C. 系统维护人员需求：}
\leaditem{快速定位}{通过协议返回的详细状态码和日志，能快速定位硬件或网络问题。}
\leaditem{远程诊断与升级}{支持远程查询终端参数、日志，并支持固件远程升级（FOTA）。}
\leaditem{安装调试简便}{协议应支持设备的自注册、自动网络接入，降低现场配置复杂度。}

\subsection{设计需求（协议非功能性需求）}
基于用户需求，推导出协议本身必须满足的设计需求：

\leaditem{可靠性}{在校园复杂的无线（如NB-IoT、LoRa、Wi-Fi）或有线环境下，保证数据传输的准确、不丢失、不重复。关键指令（如扣费、关阀）需具备确认机制。}
\leaditem{实时性}{状态变化（如开始用水、结束用水）应在数秒内上报，控制指令应在可接受延时内（如<5秒）得到执行。}
\leaditem{安全性}{}
\subleaditem{数据保密性}{对关键用户数据（如用户ID、交易记录）进行加密传输。}
\subleaditem{数据完整性}{使用报文校验（如CRC16）和数字签名，防止数据在传输中被篡改。}
\subleaditem{身份认证}{终端与服务器之间需进行双向认证，防止非法设备接入或伪服务器攻击。}
\subleaditem{防重放攻击}{重要交易指令需包含序列号或时间戳，防止指令被恶意重复发送。}
\leaditem{低功耗}{对于采用电池供电的无线终端，协议应支持心跳保活、低功耗睡眠模式，仅在必要时刻通讯，以延长设备寿命。}
\leaditem{可扩展性}{协议帧结构应预留字段，支持未来新增功能（如水温、水压监测）而不影响旧版本兼容性。}
\leaditem{易解析与高效性}{报文结构应紧凑、清晰，减少冗余，以节省流量（特别是蜂窝网络）和提高解析效率。}

\section{详细协议设计}
\subsection{协议概述}
\leaditem{设计原则}{本协议采用分层设计思想，分为传输层适配和应用层协议。应用层协议独立于底层网络，可适配TCP/IP、UDP/CoAP等多种传输方式。协议设计遵循“请求-响应”与“主动上报”相结合的模式，兼顾实时性和可靠性。}

\subsection{协议框架设计}
\subsubsection{通信模型}
\begin{center}
\setstretch{1}
\begin{tikzpicture}[font=\small]
  \node[seqhead, minimum width=4.2cm, minimum height=1.3cm, align=center] (cli) at (0,0) {智能水控终端\\(Client)};
  \node[seqhead, minimum width=4.2cm, minimum height=1.3cm, align=center] (srv) at (9,0) {中心服务器\\(Server)};
  \draw[seqarrow] ([yshift=0.35cm]cli.east) -- ([yshift=0.35cm]srv.west);
  \draw[seqarrow] ([yshift=-0.35cm]srv.west) -- ([yshift=-0.35cm]cli.east);
  \node[anchor=north, align=center] at (0,-1.0) {事件触发\\心跳维持\\状态上报};
  \node[anchor=north, align=center] at (9,-1.0) {主动控制\\参数下发\\固件升级};
\end{tikzpicture}
\end{center}

\subsubsection{连接方式}
\leaditem{主连接}{TCP长连接（有线/Wi-Fi）或CoAP/UDP（NB-IoT/LoRa）}
\leaditem{保活机制}{心跳包维持连接，心跳间隔可配置（默认300秒）}
\leaditem{断线重连}{设备检测到断线后按指数退避算法重连}

\subsubsection{通信模式}
\leaditem{设备主动上报}{用水事件、故障告警、状态变化}
\leaditem{服务器请求-设备响应}{服务器查询设备状态、参数}
\leaditem{服务器指令-设备执行}{远程控制、参数配置}
\leaditem{设备请求-服务器响应}{心跳、注册、校时}

\subsection{协议报文格式}
\subsubsection{基本帧结构}
\begin{reptable}
\begin{tabularx}{\textwidth}{|*{8}{Z|}}
\hline
\tabhead{帧头 (2)} & \tabhead{版本(1)} & \tabhead{命令字(1)} & \tabhead{序列号(2)} & \tabhead{设备ID(8)} & \tabhead{长度(2)} & \tabhead{负载(N)} & \tabhead{校验(2)}\\
\hline
0xAA55 & 0x01 & Cmd & Seq & DevID & Len & Data & CRC\\
\hline
\end{tabularx}
\end{reptable}

\subsubsection{字段详细说明}
\begin{reptable}
\begin{tabularx}{\textwidth}{|>{\bfseries}c|c|Y|Y|}
\hline
\tabhead{字段名} & \tabhead{长度} & \multicolumn{1}{c|}{\tabhead{描述}} & \multicolumn{1}{c|}{\tabhead{备注}}\\
\hline
帧头 & 2字节 & 固定值 0xAA55 & 用于帧同步\\
\hline
协议版本 & 1字节 & 主版本号(高4位)+次版本号(低4位) & 版本1.2 \ensuremath{\rightarrow} 0x12\\
\hline
命令字 & 1字节 & 指令类型标识 & 见3.4节命令字定义\\
\hline
序列号 & 2字节 & 消息序号，从0递增，循环使用 & 用于请求-响应匹配和防重放\\
\hline
设备ID & 8字节 & 设备唯一标识符 & 生产时烧录，格式: 厂商(2)+型号(2)+序号(4)\\
\hline
负载长度 & 2字节 & 负载数据的字节长度 & 大端字节序，范围0-1024\\
\hline
负载数据 & N字节 & 具体业务数据 & 结构随命令字变化\\
\hline
CRC16校验 & 2字节 & 对整个帧（除校验字段）的CRC16校验 & 多项式: 0x8005\\
\hline
\end{tabularx}
\end{reptable}

\subsection{命令字定义}
\subsubsection{基础命令集（0x00-0x3F）}
\begin{reptable}
\begin{tabularx}{\textwidth}{|>{\bfseries}c|l|l|Y|}
\hline
\tabhead{命令字} & \multicolumn{1}{c|}{\tabhead{名称}} & \multicolumn{1}{c|}{\tabhead{方向}} & \multicolumn{1}{c|}{\tabhead{说明}}\\
\hline
0x01 & 心跳上报 & \devtosrv & 设备在线状态报告\\
\hline
0x02 & 心跳响应 & \srvtodev & 服务器确认收到心跳\\
\hline
0x03 & 设备注册 & \devtosrv & 设备首次上线注册\\
\hline
0x04 & 注册响应 & \srvtodev & 注册结果，下发设备密钥\\
\hline
0x05 & 时间同步请求 & \devtosrv & 请求服务器时间\\
\hline
0x06 & 时间同步响应 & \srvtodev & 返回服务器时间戳\\
\hline
\end{tabularx}
\end{reptable}

\subsubsection{用水控制命令集（0x10-0x1F）}
\begin{reptable}
\begin{tabularx}{\textwidth}{|>{\bfseries}c|l|l|Y|}
\hline
\tabhead{命令字} & \multicolumn{1}{c|}{\tabhead{名称}} & \multicolumn{1}{c|}{\tabhead{方向}} & \multicolumn{1}{c|}{\tabhead{说明}}\\
\hline
0x10 & 用水开始事件 & \devtosrv & 用户刷卡/扫码成功，开始用水\\
\hline
0x11 & 用水结束事件 & \devtosrv & 用户停止用水或超时停止\\
\hline
0x12 & 实时用水上报 & \devtosrv & 用水过程中定时上报累积水量\\
\hline
0x13 & 远程开阀指令 & \srvtodev & 管理员远程开阀\\
\hline
0x14 & 远程关阀指令 & \srvtodev & 管理员远程关阀\\
\hline
0x15 & 开关阀响应 & \devtosrv & 执行开关阀操作的结果\\
\hline
\end{tabularx}
\end{reptable}

\subsubsection{状态查询命令集（0x20-0x2F）}
\begin{reptable}
\begin{tabularx}{\textwidth}{|>{\bfseries}c|l|l|Y|}
\hline
\tabhead{命令字} & \multicolumn{1}{c|}{\tabhead{名称}} & \multicolumn{1}{c|}{\tabhead{方向}} & \multicolumn{1}{c|}{\tabhead{说明}}\\
\hline
0x20 & 设备状态查询 & \srvtodev & 查询设备当前状态\\
\hline
0x21 & 设备状态上报 & \devtosrv & 上报设备状态信息\\
\hline
0x22 & 累计数据查询 & \srvtodev & 查询设备累计用水量等信息\\
\hline
0x23 & 累计数据上报 & \devtosrv & 上报累计数据\\
\hline
0x24 & 实时数据查询 & \srvtodev & 查询设备实时数据（温度、流量等）\\
\hline
0x25 & 实时数据上报 & \devtosrv & 上报实时数据\\
\hline
\end{tabularx}
\end{reptable}

\subsubsection{参数配置命令集（0x30-0x3F）}
\begin{reptable}
\begin{tabularx}{\textwidth}{|>{\bfseries}c|l|l|Y|}
\hline
\tabhead{命令字} & \multicolumn{1}{c|}{\tabhead{名称}} & \multicolumn{1}{c|}{\tabhead{方向}} & \multicolumn{1}{c|}{\tabhead{说明}}\\
\hline
0x30 & 参数读取请求 & \srvtodev & 读取设备参数\\
\hline
0x31 & 参数读取响应 & \devtosrv & 返回设备参数\\
\hline
0x32 & 参数设置指令 & \srvtodev & 设置设备参数\\
\hline
0x33 & 参数设置响应 & \devtosrv & 参数设置结果\\
\hline
0x34 & 费率参数下发 & \srvtodev & 下发费率表\\
\hline
0x35 & 费率参数确认 & \devtosrv & 费率接收确认\\
\hline
\end{tabularx}
\end{reptable}

\subsubsection{故障告警命令集（0x40-0x4F）}
\begin{reptable}
\begin{tabularx}{\textwidth}{|>{\bfseries}c|l|l|Y|}
\hline
\tabhead{命令字} & \multicolumn{1}{c|}{\tabhead{名称}} & \multicolumn{1}{c|}{\tabhead{方向}} & \multicolumn{1}{c|}{\tabhead{说明}}\\
\hline
0x40 & 故障实时上报 & \devtosrv & 故障发生时立即上报\\
\hline
0x41 & 故障确认指令 & \srvtodev & 服务器确认收到故障\\
\hline
0x42 & 故障恢复上报 & \devtosrv & 故障恢复时上报\\
\hline
0x43 & 告警查询指令 & \srvtodev & 查询设备告警记录\\
\hline
0x44 & 告警记录上报 & \devtosrv & 上报告警记录\\
\hline
\end{tabularx}
\end{reptable}

\subsubsection{固件升级命令集（0x50-0x5F）}
\begin{reptable}
\begin{tabularx}{\textwidth}{|>{\bfseries}c|l|l|Y|}
\hline
\tabhead{命令字} & \multicolumn{1}{c|}{\tabhead{名称}} & \multicolumn{1}{c|}{\tabhead{方向}} & \multicolumn{1}{c|}{\tabhead{说明}}\\
\hline
0x50 & 升级通知指令 & \srvtodev & 通知设备有新固件\\
\hline
0x51 & 升级请求 & \devtosrv & 设备请求升级\\
\hline
0x52 & 升级数据包 & \srvtodev & 下发固件数据包\\
\hline
0x53 & 升级确认 & \devtosrv & 数据包接收确认\\
\hline
0x54 & 升级结果上报 & \devtosrv & 升级成功/失败结果\\
\hline
\end{tabularx}
\end{reptable}

\subsection{关键业务流程}
\subsubsection{设备注册流程}
\begin{seqdiagram}
  \seqtitles{设备}{服务器}
  \seqlines{4.6}
  \seqsend{1.1}{注册请求(0x03)}
  \seqnote{1.8}{1. 验证设备ID合法性}
  \seqnote{2.5}{2. 生成设备专属会话密钥}
  \seqrecv{3.2}{注册响应(0x04)}
  \seqnote{3.2}{3. 返回注册结果+密钥}
  \seqnote{3.9}{4. 设备存储密钥}
\end{seqdiagram}

\subsubsection{用水控制流程}
\begin{seqdiagram}
  \seqlines{13.2}
  \seqact{0.6}{用户刷卡/扫码}
  \seqact{1.4}{设备验证卡片/二维码}
  \seqact{2.2}{发送用水开始事件(0x10)}
  \seqsend{3.0}{}
  \seqnote{3.0}{1. 记录开始时间}
  \seqnote{3.7}{2. 验证用户余额}
  \seqrecv{4.4}{}
  \seqnote{4.4}{3. 返回授权结果}
  \seqact{5.2}{开阀放水}
  \seqact{6.0}{每30秒上报实时用水(0x12)}
  \seqsend{6.8}{}
  \seqnote{6.8}{记录实时水量}
  \seqact{7.6}{用户停止用水/超时}
  \seqact{8.4}{关阀停止}
  \seqact{9.2}{发送用水结束事件(0x11)}
  \seqsend{10.0}{}
  \seqnote{10.0}{1. 计算消费金额}
  \seqnote{10.7}{2. 扣除用户余额}
  \seqnote{11.4}{3. 生成消费记录}
  \seqrecv{12.1}{}
  \seqnote{12.1}{4. 返回确认}
\end{seqdiagram}

\subsubsection{远程控制流程}
\begin{seqdiagram}
  \seqlines{5.8}
  \seqact{0.6}{后台管理员}
  \seqact{1.4}{选择设备，点击“远程开阀”}
  \seqact{2.2}{发送开阀指令(0x13)}
  \seqsend{3.0}{}
  \seqnote{3.0}{验证管理员权限}
  \seqrecv{3.8}{}
  \seqnote{3.8}{执行开阀操作}
  \seqact{4.6}{发送开关阀响应(0x15)}
  \seqsend{5.4}{}
  \seqnote{5.4}{记录操作日志}
\end{seqdiagram}

\subsubsection{故障上报流程}
\begin{seqdiagram}
  \seqlines{6.4}
  \seqact{0.6}{设备检测到故障}
  \seqact{1.4}{立即发送故障上报(0x40)}
  \seqsend{2.2}{}
  \seqnote{2.2}{1. 生成故障工单}
  \seqnote{2.9}{2. 推送告警信息}
  \seqrecv{3.6}{}
  \seqnote{3.6}{3. 发送故障确认(0x41)}
  \seqact{4.4}{故障恢复}
  \seqact{5.2}{发送故障恢复上报(0x42)}
  \seqsend{6.0}{}
  \seqnote{6.0}{关闭故障工单}
\end{seqdiagram}

\subsection{数据格式定义}
\subsubsection{用水开始事件（0x10）负载格式}
\begin{reptable}
\begin{tabularx}{\textwidth}{|*{6}{Z|}}
\hline
用户ID(8) & 卡片类型 & 余额(4) & 费率ID(2) & 预留(2) & 时间戳(4)\\
\hline
\end{tabularx}
\end{reptable}

\subsubsection{用水结束事件（0x11）负载格式}
\begin{reptable}
\begin{tabularx}{\textwidth}{|*{7}{Z|}}
\hline
\fld{用户ID}{8} & \fld{结束水量}{4} & \fld{用水时长}{4} & \fld{消费金额}{4} & \fld{原因}{1} & \fld{预留}{3} & \fld{时间戳}{4}\\
\hline
\end{tabularx}
\end{reptable}

\subsubsection{设备状态（0x21）负载格式}
\begin{reptable}
\begin{tabularx}{\textwidth}{|*{8}{Z|}}
\hline
\fld{阀门状态}{1} & \fld{水温}{2} & \fld{水压}{2} & \fld{信号强度}{1} & \fld{电压}{2} & \fld{故障码}{2} & \fld{工作模式}{1} & \fld{预留}{2}\\
\hline
\end{tabularx}
\end{reptable}

\subsubsection{故障上报（0x40）负载格式}
\begin{reptable}
\begin{tabularx}{\textwidth}{|*{7}{Z|}}
\hline
\fld{故障代码}{2} & \fld{故障级别}{1} & \fld{故障数值}{4} & \fld{相关参数1}{4} & \fld{相关参数2}{4} & \fld{预留}{2} & \fld{时间戳}{4}\\
\hline
\end{tabularx}
\end{reptable}

\subsection{安全设计}
\subsubsection{加密方案}
\leaditem{传输层}{强制使用TLS 1.2+（TCP连接）或DTLS（UDP连接）}
\leaditem{应用层}{敏感数据使用AES-128-GCM加密}
\leaditem{密钥管理}{设备出厂预置根密钥，注册时更新为会话密钥}

\subsubsection{认证机制}
\leaditem{设备认证}{基于预置证书/密钥的TLS双向认证}
\leaditem{消息认证}{每条消息携带基于会话密钥的HMAC-SHA256}
\leaditem{用户认证}{刷卡/扫码时与后台系统交互认证}

\subsubsection{防重放攻击}
\noindent 每个消息包含递增序列号（2字节，循环使用）\par
\noindent 服务器维护设备最近N个序列号的滑动窗口\par
\noindent 时间戳校验，允许±5分钟的时间误差\par

\subsubsection{数据完整性}
\leaditem{帧级}{CRC16校验}
\leaditem{消息级}{HMAC消息认证码}
\leaditem{应用级}{关键业务数据包含数字签名}

\subsection{错误处理}
\subsubsection{错误码定义}
\begin{reptable}
\setlength{\tabcolsep}{14pt}
\begin{tabular}{|c|l|l|}
\hline
\tabhead{错误码} & \multicolumn{1}{c|}{\tabhead{含义}} & \multicolumn{1}{c|}{\tabhead{处理建议}}\\
\hline
0x00 & 成功 & 正常处理\\
\hline
0x01 & 协议版本不支持 & 升级设备固件\\
\hline
0x02 & 命令字不支持 & 检查协议版本\\
\hline
0x03 & 数据格式错误 & 检查数据格式\\
\hline
0x04 & 参数错误 & 检查参数范围\\
\hline
0x05 & 认证失败 & 重新注册或检查密钥\\
\hline
0x06 & 权限不足 & 检查操作权限\\
\hline
0x07 & 设备忙 & 稍后重试\\
\hline
0x08 & 内存不足 & 清理设备数据\\
\hline
0x09 & 网络异常 & 检查网络连接\\
\hline
0x0A & 服务器忙 & 稍后重试\\
\hline
0x0B-0x7F & 预留 & \\
\hline
\end{tabular}
\end{reptable}

\subsubsection{重传机制}
\leaditem{关键指令（扣费、控制）}{最多重试3次，间隔2秒、4秒、8秒}
\leaditem{普通数据}{最多重试1次}
\leaditem{心跳包}{不重传，超时直接重连}

\subsection{性能指标}
\subsubsection{通信性能}
\begin{reptable}
\setlength{\tabcolsep}{14pt}
\begin{tabular}{|c|c|c|}
\hline
\tabhead{指标} & \tabhead{目标值} & \tabhead{备注}\\
\hline
事件上报延迟 & <3秒 & 用水开始/结束\\
\hline
控制指令响应 & <5秒 & 远程开关阀\\
\hline
状态查询响应 & <2秒 & 实时状态查询\\
\hline
心跳间隔 & 300秒 & 可配置\\
\hline
最大包大小 & 1024字节 & 固件升级除外\\
\hline
\end{tabular}
\end{reptable}

\subsubsection{网络流量估算}
\begin{reptable}
\setlength{\tabcolsep}{14pt}
\begin{tabular}{|c|c|c|c|}
\hline
\tabhead{场景} & \tabhead{上行流量} & \tabhead{下行流量} & \tabhead{频率}\\
\hline
心跳包 & 30字节 & 20字节 & 每5分钟\\
\hline
用水开始事件 & 50字节 & 20字节 & 每次用水\\
\hline
用水结束事件 & 60字节 & 20字节 & 每次用水\\
\hline
实时用水上报 & 40字节 & 0字节 & 每30秒（用水中）\\
\hline
月均用量估算 & \textasciitilde 5MB & \textasciitilde 2MB & 按每天50次用水\\
\hline
\end{tabular}
\end{reptable}

\subsection{扩展性设计}
\subsubsection{版本兼容}
\leaditem{主版本号不同}{不兼容，拒绝服务}
\leaditem{次版本号不同}{向后兼容，忽略未知字段}
\noindent 扩展字段使用TLV格式，旧版本可跳过未知类型\par

\subsubsection{功能扩展预留}
\leaditem{负载预留字段}{每个消息负载保留4字节扩展位}
\leaditem{命令字空间}{0x80-0xFF为厂商自定义命令}
\leaditem{设备能力协商}{注册时上报支持的功能列表}

\subsubsection{未来可扩展功能}
\noindent 水温自适应调节（需增加温度传感器）\par
\noindent 用水习惯分析（需增加数据分析字段）\par
\noindent 语音提示功能（需增加语音控制命令）\par
\noindent 多人用水模式（需增加多用户管理）\par

\unnumberedsection{附\quad 录}

\appitem{数据类型定义}
\begin{lstlisting}[language=C]
// 基本数据类型
typedef uint8_t  BYTE;       // 1字节无符号
typedef uint16_t WORD;       // 2字节无符号
typedef uint32_t DWORD;      // 4字节无符号
typedef uint64_t QWORD;      // 8字节无符号
typedef int32_t  LONG;       // 4字节有符号

// 时间类型：Unix时间戳（秒）
typedef DWORD TIMESTAMP;

// 金额类型：单位分（1元=100分）
typedef LONG MONEY;
\end{lstlisting}

\appitem{故障码表（部分）}
\noindent 0x0001: 阀门故障 - 开阀失败\\
0x0002: 阀门故障 - 关阀失败\\
0x0003: 流量计故障 - 无脉冲信号\\
0x0004: 温度传感器故障\\
0x0005: 水压传感器故障\\
0x0006: 存储器故障\\
0x0007: 电源电压过低\\
0x0008: 通讯模块故障\\
0x0009: 余额不足警告\\
0x000A: 非法卡警告\par

\appitem{协议示例}
\begin{lstlisting}[language=Python]
# 用水开始事件示例（0x10）
frame = {
    'header': 0xAA55,
    'version': 0x01,
    'cmd': 0x10,
    'seq': 0x0102,
    'dev_id': 0x1234567890ABCDEF,
    'length': 24,
    'payload': {
        'user_id': 0x1001,
        'card_type': 1,      # 1:校园卡，2:微信，3:支付宝
        'balance': 5000,     # 50.00元
        'rate_id': 1,
        'reserved': 0,
        'timestamp': 1672502400
    },
    'crc': 0x1234
}
\end{lstlisting}

\appitem{CRC-16计算示例}
\blocklabel{示例1：心跳包（命令字 0x01）}
\noindent 场景：设备ID为 0x1122334455667788 的设备发送心跳包\par

\blocklabel{步骤1：构建数据包（不含CRC）}
\noindent 原始数据（16进制表示）：
\begin{lstlisting}[language={}]
AA 55 01 01 01 02 11 22 33 44 55 66 77 88 00 00
\end{lstlisting}
\noindent 字段解析：
\begin{center}
\setstretch{1}
\begin{tabular}{@{}l@{\quad-\quad}l@{}}
\ttfamily AA 55 & 帧头 (0xAA55)\\
\ttfamily 01 & 版本 (V1.1)\\
\ttfamily 01 & 命令字 (0x01 = 心跳)\\
\ttfamily 01 02 & 序列号 (0x0102)\\
\ttfamily 11 22 33 44 55 66 77 88 & 设备ID\\
\ttfamily 00 00 & 负载长度 (0字节，无负载)\\
\end{tabular}
\end{center}

\blocklabel{步骤2：计算CRC}
\begin{lstlisting}[language=Python]
data = bytes.fromhex('AA550101010211223344556677880000')
crc = CRC16_MODBUS.calculate(data)
print(f"CRC16校验值: 0x{crc:04X}")  # 输出: 0xB9A5
\end{lstlisting}

\blocklabel{步骤3：添加CRC得到完整数据包}
\noindent 完整数据包（16进制）：
\begin{lstlisting}[language={}]
AA 55 01 01 01 02 11 22 33 44 55 66 77 88 00 00 B9 A5
\end{lstlisting}
\noindent 字节位置说明：
\begin{center}
\setstretch{1}
\begin{tabular}{@{}l@{:\ }l@{\quad}l@{}}
\ttfamily 00-01 & 帧头 & \ttfamily 0xAA55\\
\ttfamily 02 & 版本 & \ttfamily 0x01\\
\ttfamily 03 & 命令字 & \ttfamily 0x01\\
\ttfamily 04-05 & 序列号 & \ttfamily 0x0102\\
\ttfamily 06-13 & 设备ID & \ttfamily 0x1122334455667788\\
\ttfamily 14-15 & 长度 & \ttfamily 0x0000\\
\ttfamily 16-17 & CRC16 & \ttfamily 0xB9A5\\
\end{tabular}
\end{center}

\blocklabel{步骤4：服务器验证CRC}
\begin{lstlisting}[language=Python]
# 服务器端验证
full_packet = bytes.fromhex('AA550101010211223344556677880000B9A5')
is_valid = CRC16_MODBUS.verify(full_packet)
if is_valid:
    print("CRC验证通过")
else:
    print("CRC验证失败，数据可能被篡改")
\end{lstlisting}

\end{document}
